% ==============================================================================
% Frontmatter: Title, Abstract, Keywords
% ==============================================================================
\begin{frontmatter}
\title{\textbf{Frequent Losers as Cover Bidders in Public Procurement\tnoteref{tfoot}}}
\tnotetext[tfoot]{\emph{This version: March 2026}}
\author[INSPER]{Darcio Genicolo-Martins} \ead{darciogm1@insper.edu.br}
\author[INSPER]{Paulo Furquim de Azevedo}
\address[INSPER]{INSPER --- Institute of Education and Research, S\~{a}o Paulo, Brazil}

\begin{abstract}
We propose \textit{frequent losers} (FL)---firms that systematically participate in and lose public procurement tenders---as a novel firm-level screening marker for cover bidding in cartel arrangements. Using 4.5 million tender-items and 40 million bids from Brazil's BEC electronic procurement platform (2009--2019), we establish three contributions. First, we show that FL presence in a tender is associated with 6--9\% higher winning prices across multiple specifications, robustness checks, and matching estimators, and we validate this association externally against CADE (Brazilian Competition Authority) cartel convictions: 193 FL firms co-participate in tenders with convicted cartelists. Second, we empirically distinguish between complementary cover bidding (Regime 1), where FL bids are drawn independently above the reference price, and coordinated cover bidding (Regime 2), where FL bids cluster near the winning price. Our data are consistent with Regime~1: FL bid dispersion is significantly higher than non-FL dispersion, and 91\% of FL bids exceed the winning price with a median spread of 55\%. Third, Bajari \& Ye (2003) tests reject both exchangeability and conditional independence for FL bid residuals, providing direct statistical evidence that FL firms bid differently and in a correlated manner within tenders. FL-present tenders also exhibit more total participants but at least as many genuine competitors, consistent with cartels strategically deploying cover bidders in tenders with greater legitimate competition. Back-of-envelope welfare calculations suggest the FL-associated markup costs S\~{a}o Paulo R\$240--1,250 million over the sample period, or approximately 6.6\% of winning prices in FL-present tenders.
\end{abstract}

\end{frontmatter}
\textbf{Keywords:} bid rigging, cover bidding, public procurement, cartel screening, frequent losers\\
\textbf{JEL No.:} D44, D73, H57, K21, L41
